% Informal single-component picture; technical qualifications stay in the notes.
\begin{frame}{Factorizable Attacks}
  \centering
  \setbeamertemplate{theorems}[normal font]
  \setbeamercolor{block title}{fg=white,bg=CVQCInk}
  \setbeamercolor{block body}{fg=CVQCInk,bg=CVQCFog}
  \setbeamercolor{itemize item}{fg=CVQCTeal}
  \begin{tikzpicture}
    \path[use as bounding box] (-6.6,-3.15) rectangle (6.6,3.05);

    \node[inner sep=0pt] at (0,1.05) {%
      \begin{minipage}{11.4cm}
        \raggedright
        \begin{definition}[Informal]
          \par\vspace{0.5em}\noindent
          An attack is \textbf{factorizable} if it splits across the pad:
          {\large
          \[
            \textcolor{CVQCCoral}{A_{\mathrm{ct}}}\,
            \operatorname{OTP}_k\,
            \textcolor{CVQCCoral}{A_{\mathrm{pt}}}\ket{\psi}.
          \]
          }
          \noindent Both attacks are independent of the pad key $k$.
        \end{definition}
      \end{minipage}%
    };

    \node[anchor=north west,inner sep=0pt,text=CVQCInk] at (-5.5,-0.95) {%
      \begin{minipage}{11cm}
        \raggedright
        \begin{itemize}
          \setlength{\itemsep}{0.75em}
          \item<2-> Subsumes blindness.
          \item<3-> $\textcolor{CVQCCoral}{A_{\mathrm{ct}}}$:
            Statistical analogue of an ``up to isometry'' guarantee.
        \end{itemize}
      \end{minipage}%
    };
  \end{tikzpicture}

  \note<1>{
    \footnotesize
    The previous output had BQ outside the pad and X-a inside. Name those
    pieces ciphertext-level and plaintext-level attacks. Psi here denotes
    the intended plaintext, including the honest target computation.
    The displayed product is one component in the attack decomposition;
    the complete state sums components and averages uniform reference keys.
    Private auxiliary registers and their initially key-independent state
    are implicit. The outside attack may act globally, including on the
    prover's memory. Both attack operators are independent of the one-time
    pad key, although they may depend on public instructions and analysis
    indices. The verifier keeps the prescribed key bookkeeping. In the
    MBQC characterization the plaintext modification is unitary; later
    distillation uses the formal copywise attack-map definition, with its
    subnormalization and aggregate trace bounds. Source: Technical Overview,
    Factorization of attacks; Real-Copy Protocol, output characterization.
  }
  \note<2>{
    \footnotesize
    Explain the blindness point in the setting of our complete MBQC
    characterization. The reference pad is uniform, the plaintext
    modifications are unitary, and the residual ciphertext operators do
    not depend on that key. Averaging the pad hides the modified plaintext;
    the ciphertext attack then acts on a hidden state. The proof tracks
    the transcript, memory, and normalization jointly. This is stronger
    structural information than a game asserting privacy alone. An isolated
    arbitrary trace-decreasing plaintext operator is not by itself a
    general blindness theorem. The formal statement and its normalization
    conditions supply the guarantee. This structure is what prevents the
    later distillation input from carrying the secret key in private garbage.
  }
  \note<3>{
    \footnotesize
    For the RSP analogy, take $\operatorname{OTP}_k=X^xT^t$ and
    $\ket{\psi}=\ket{+}$. If $A_{\mathrm{pt}}=I$, the intended state is
    $X^xT^t\ket{+}=\ket{+_\theta}$ up to global phase, with
    $\theta=(-1)^x t\pi/4$ modulo $2\pi$. Only the key-independent outside
    operation remains. This resembles GV19's accepted-state representation:
    an angle-independent isometry acting on the ideal RSP state and an
    auxiliary state (Section 4, Eq. 32,
    \href{https://arxiv.org/pdf/1904.06320}{arXiv:1904.06320}).
    Our MBQC characterization is statistical, even against an unbounded
    prover. An individual $A_{\mathrm{ct}}$ can decrease norm, so it need
    not itself be an isometry. The complete trace-preserving operation,
    including any rejection flag, admits an isometric dilation when its
    environment is retained. This is an analogy between state forms;
    it does not by itself establish the same composable RSP functionality.
    Next, explain why distilling these attacked ciphertexts is a challenge.
  }
\end{frame}
